\documentclass[10pt]{article}

\usepackage[utf8]{inputenc}

\usepackage[T1]{fontenc}

\usepackage{amsmath}

\usepackage{amsfonts}

\usepackage{amssymb}

\usepackage[version=4]{mhchem}

\usepackage{stmaryrd}

\usepackage{graphicx}

\usepackage[export]{adjustbox}

\graphicspath{ {./images/} }

\usepackage{bbold}



\begin{document}

кальных П. т. описывает предельное распределение плотностей распределения сумм.



\begin{enumerate}

  \setcounter{enumi}{3}

  \item П.т. в их классич. постановке описывают поведение отдельной суммы $s_{n}$ с возрастанием номера $n$. Достаточно общие П. т. для вероятностей событий, зависящих сразу от нескольких сумм, получены впервые А. Н. Колмогоровым (1931). Так, напр., из его результатов следует, что при весьма пироких условиях вероятность неравенства

\end{enumerate}



\[

\max _{1 \leqslant k \leqslant n}\left|s_{k}\right|<z B_{n}

\]



имеет пределом величину



\[

\frac{4}{\pi} \sum_{k=0}^{\infty}(-1)^{k} \frac{1}{2 k+1} e^{-(2 k+1)^{2} z^{2} / 8 \pi^{2}}, \quad z>0 .

\]



Наиболее общий способ доказательства подобных П. т. - предельный переход от дискретных процессов к непрерывным.



\begin{enumerate}

  \setcounter{enumi}{4}

  \item Перечисленные выше П. т. относятся к суммам случайных величин. Примером П. т. иного рода могут служить П.т. для членов вариационного ряда. Эти П. т. подробно изучены Б. В. Гнеденко, Н. В. Смирновым и др.



  \item Наконец, к П. т. относят также и теоремы, устанавливающие свойства последовательностей случайных величин, имеющие место с вероятностью, равной единице (см., напр., Больших чисел усиленный закон, Повторного логарифла закон).



\end{enumerate}



О методах доказательства П.т. см. в статьях $X a^{-}$ рактеристическая бункция, Распределений сходиMocmb.



Лum.: [1] Гне ден ко Б. В., Колмогоров А. Н., Предельные распределения для сумм независимых случайных величин, М.- Л., 1949; [2] И б р а г и м о в И. А., Л и нн и к Ю. В., Независимые и стационарно связанные величины, М., 1965; [3]' П е т р о в В. В., Суммы независимых случайных



\includegraphics[max width=\textwidth, center]{2023_07_08_c578af8fe2826f436851g-1(1)}

Ю. А., Теория вероптос, Ю. В. Прохоров.



ПРЕДЕЛЫНЫИ КОНУС В Ы П У К Л о Й П о В е р хн о с т и $S$ - поверхность $P(S)$ конуса, образованного полупрямыми, исходящими из нек-рой точки $O \in S$ и принадлежащими выпуклому телу, ограниченному $S$. П. к. определен однозначно с точностью до параллельного переноса, зависящего от выбора точки $O$. Понятие П. к. определяется и для нек-рых классов невыпуклых поверхностей, напр. для т. н. сферически однолистных седловых поверхностей.



Лит.: [1] П о г о р е ло в А. В., Внешняя геометрия выпуклых поверхностей, М. 1969 . М. И. Войцеховский.



ПРЕДЕЛЬНЫД ЦИКЛ - замкнутая траектория в фазовом пространстве автономной системы обыкновенных дифференциальных уравнений, к-рая является $\alpha$ - или $\omega$-предельным множеством (см. IIредельное множество траектории) хотя бы для одной другой траектории этой системы. П. ц. наз. о р б и т а л ь н о у с то й ч и вы м, или уст о й ч и вым, если для всякого $\varepsilon>0$ найдется $\delta>0$ такое, что все траектории, начинающиеся в $\delta$-окрестности П. ц. при $t=0$, не вы-



\includegraphics[max width=\textwidth, center]{2023_07_08_c578af8fe2826f436851g-1}

вует периодич. решение системы, отличное от постоянного. Если оно устойчиво по Ляпунову, то П. Ц. устойчив. Для того чтобы периодич. решению соответствовал устойчивый П. ц., достаточно, чтобы модули всех его мультипликаторов, кроме одного, были меньше единицы (см. Орбитальнал устойчивость, Андронова - Bитта теорема). С физич. точки зрения П. ц. соответствует периодич. режиму, или автоколебанию, системы (см. [2]).



Пусть автономная система



\[

\dot{x}=f(x)

\]



руемое многообразие, напр. $V^{n}=\mathbb{R} n$, имеет замкнутую траекторию $\Gamma$. Проведем гиперповерхность $\pi$, пересекающую $\Gamma$ трансверсально в точке $p$. Тогда любая траектория системы, начинающаяся при $t=0$ в точке $c \in V \subset \pi$, где $V$ - достаточно малая окрестность точки $p$, при увеличении $t$ снова пересечет $\pi$ в точке $T(c)$. Диффеоморфизм $T: V \rightarrow T(V)$ имеет неподвижную точку $p$ и наз. последования отображением. Его свойства определяют поведение траекторий системы в окрестности Г. П. ц. в отличие от произвольной замкнутой траектории всегда определяет отличное от тождественного отображение последования. Если $p$ - седловая точка диффеоморфизма $T$, то П. ц. наз. П. ц. с е д л о в о г о т и п а. Систөма, имеющая П. ц. седлового типа, может иметь гомоклинич. кривую, т. е. траекторию, для к-рой П. Д. является одновременно как $\alpha-$, так и $\omega$-предельным множеством.



В случае двумерной системы $(*)\left(V^{n}=\mathbb{R}^{2}\right)$ в качестве $\pi$ берут прямую и рассматривают функцию $\rho: \rho(c)=T(c)-c$, к-рая наз. Ф у н к ци е й п о с лед о в а н и я. Кратность нуля $c=p$ функции $\rho$ наз. к р а т о о т ю П. П. Предельный цикл четной кратности наз. п о л у у с т о й ч и в ы м. П. ц. вместе с точками покоя и сепаратрисами определяют качественную картину поведения остальных траекторий (см. Пуанкаре - Бендиксона теория, а также [3], [4]). П. ц. в случае аналитич. функций $f(x)$ принадлежит к одному из типов: 1) устойчивый, 2) неустойчивый, т. е. устойчивый П. ц., если изменить направление $t$ на противоположное, 3) полуустойчивый. Напр., система



\[

\begin{gathered}

\dot{x}_{1}=-\omega x_{2}+\mu x_{1}\left(x_{1}^{2}+x_{1}^{2}-1\right)^{k}, \\

\dot{x}_{2}=\omega x_{1}+\mu x_{2}\left(x_{1}^{2}+x_{2}^{2}-1\right)^{k},

\end{gathered}

\]



где $\omega \neq 0, \quad \mu=$ const, $k \in \mathbb{N}$, имеет при $\mu<0(\mu>0)$ и $k$ неqетном устойчивый (неустойчивый), а при $k$ четном - полуустойчивый П.ц. кратности $k$; во всех этих случаях П.ц. является окружность $x_{1}^{2}+x_{2}^{2}=1$, т. е. траектория решения



\[

x_{1}=\cos \left(\omega t+\varphi_{0}\right), x_{2}=\sin \left(\omega t+\varphi_{0}\right) .

\]



Если система (*) задана на односвязной области $U \subset \mathbb{R}^{2}$, то П. ц. окружает по крайней мере одну точку покоя этой системы.



Для разыскания П. ц. в системе 2-го порядка применяется метод, основанный на следующем утверждении: если векторное поле $f(x)$ направлено вовнутрь (вовне) кольцеобразной области $G$ и в $G$ нет точек покоя, то в $G$ имеется хотя бы один устойчивый (неустойчивый) П.ц. Выбор области $G$ производится из физич. соображений или результатов аналитических или численных расчетов.



Лuт.: [1] П о н т р я г и в Л. С., Обыкновенные дифференциальные уравнения, 4 изд., М., 1974; [2] А н д р о н о в А. А., В и т т А. А., Х а й к и н С. Э., Теория колебаний, 2 изд., Го р дон И. А., Ма й е р А. Г., Качественная теория динамических систем второго порядка, М., 1966; [4] и х ж е, Теория бифуркаций пинамических систем на плоскости, М., 1967; М о и 0 e в нелинейной механики, 2 изд., М., і981. Л. А. Черкас. ПРЕДИКАТ - функция, значениями к-рой являются высказывания об $n$-ках объектов, представляющих значения аргументов; при $n=1$ П. наз. «свойством», при $n>1$ - «отнотением», единичные высказывания могут рассматриваться как нульместные П.



Чтобы задать $n$-местный предикат $P\left(x_{1}, \ldots, x_{n}\right)$, следует указать множества $D_{1}, \ldots, D_{n}$ - области изменения п р е д м е т н х п е р е ме н вы $x_{1}, \ldots, x_{n}$, причем чаще всего рассматривают случай $D_{1}=D_{2}=\ldots=D_{n}$. С теоретико-множественной точки зрения П. определяется заданием подмвожества $M$





\end{document}